\documentclass[10pt,a4paper]{amsart}
\usepackage[margin=19mm]{geometry}
\usepackage{amsmath,amssymb,mathtools}
\usepackage[scr=rsfs,cal=euler]{mathalfa}
\usepackage{xfrac}
\usepackage[unicode,hidelinks,bookmarks=false]{hyperref}
\newcommand{\ie}{i.e.\ }
\newcommand{\cf}{cf.\ }
\newcommand{\resp}{respectively\ }
\newcommand{\Subsection}{\subsection}
\providecommand{\nobreakdash}{\nobreak\hbox{-}}
\setlength{\emergencystretch}{2em}
\multlinegap0pt
\usepackage{amssymb}
\usepackage[shortlabels]{enumitem}
\newtheorem{proposition}{Proposition}
\title{A 50-Vertex Cubic Counterexample to the\\Domination-versus-Edge-Domination Conjecture}
\author{Koyar Afrasyab}
\date{Source: arXiv:2609.10783v1 (CC BY 4.0)}
\begin{document}
\maketitle

\noindent\emph{Source and licence.} A 50-Vertex Cubic Counterexample to the\\Domination-versus-Edge-Domination Conjecture. Version arXiv:2609.10783v1 (CC BY 4.0), distributed by arXiv under the Creative Commons Attribution 4.0 International licence, http://creativecommons.org/licenses/by/4.0/, as recorded in the arXiv licence metadata for this exact version and shown on the abstract page https://arxiv.org/abs/2609.10783v1. Full licence notice: \texttt{licenses/arxiv-2609.10783-CC-BY-4.0.txt}.\par\smallskip

\noindent\textit{Editorial scope.} Counted unit: the article's proposition prop:edge (source0.tex lines 103--105). The article's complete original proof of it is delivered with it (source0.tex lines 107--115, one \texttt{proof} environment of 9 lines). No mathematics is written, repaired or re-derived: the statement and the proof are the article's own lines, in the article's own notation, and no retained line is substituted or re-flowed.  Retained uncounted as the source-local context the proof needs: lines 72--72.  Nothing was omitted from any retained span, so the package carries no omission record.  No retained line contains a citation, so the wrapper carries no bibliography and no bibliography entry.  No retained line contains a figure environment, an image-inclusion command or drawing code, so no image is delivered and the package declares no asset dependency.  Editorial formatting only, in addition: an AMS \texttt{amsart} preamble that restates the article's own declarations and macros used by the retained lines, namely the environments \texttt{proposition} on separate counters, together with the standard packages those lines need.  The retained environments print in this document as Proposition 1 in order of appearance, whereas the article prints the same items with the numbering of its own full document.  The complete native source of the article as distributed in the arXiv e-print for this exact version is bundled unchanged as \texttt{sources/209949/source0.tex}.
\section{The graph}\label{sec:graph}

\begin{proposition}\label{prop:edge}
For the graph \(G\) of Section~\ref{sec:graph}, \(\gamma_e(G)=15\).
\end{proposition}

\begin{proof}
Let
\[
M=\{v_j^-v_j^+:1\leq j\leq15\}.
\]
The 20 vertices not saturated by \(M\) are precisely the clause vertices, and no two clause vertices are adjacent.  Thus \(M\) is maximal, so \(\gamma_e(G)\leq15\).

Conversely, in a cubic graph an edge is adjacent to at most four other edges, two at each endpoint.  Hence a single edge can dominate at most five edges including itself.  Since \(G\) has 75 edges, every edge-dominating set has size at least \(75/5=15\).  Therefore \(\gamma_e(G)=15\).
\end{proof}

\end{document}
